% generated by GAPDoc2LaTeX from XML source (Frank Luebeck)
\documentclass[a4paper,11pt]{report}

            \usepackage{a4wide}
            \newcommand{\bbZ}{\mathbb{Z}}
        
\usepackage[top=37mm,bottom=37mm,left=27mm,right=27mm]{geometry}
\sloppy
\pagestyle{myheadings}
\usepackage{amssymb}
\usepackage[utf8]{inputenc}
\usepackage{makeidx}
\makeindex
\usepackage{color}
\definecolor{FireBrick}{rgb}{0.5812,0.0074,0.0083}
\definecolor{RoyalBlue}{rgb}{0.0236,0.0894,0.6179}
\definecolor{RoyalGreen}{rgb}{0.0236,0.6179,0.0894}
\definecolor{RoyalRed}{rgb}{0.6179,0.0236,0.0894}
\definecolor{LightBlue}{rgb}{0.8544,0.9511,1.0000}
\definecolor{Black}{rgb}{0.0,0.0,0.0}

\definecolor{linkColor}{rgb}{0.0,0.0,0.554}
\definecolor{citeColor}{rgb}{0.0,0.0,0.554}
\definecolor{fileColor}{rgb}{0.0,0.0,0.554}
\definecolor{urlColor}{rgb}{0.0,0.0,0.554}
\definecolor{promptColor}{rgb}{0.0,0.0,0.589}
\definecolor{brkpromptColor}{rgb}{0.589,0.0,0.0}
\definecolor{gapinputColor}{rgb}{0.589,0.0,0.0}
\definecolor{gapoutputColor}{rgb}{0.0,0.0,0.0}

%%  for a long time these were red and blue by default,
%%  now black, but keep variables to overwrite
\definecolor{FuncColor}{rgb}{0.0,0.0,0.0}
%% strange name because of pdflatex bug:
\definecolor{Chapter }{rgb}{0.0,0.0,0.0}
\definecolor{DarkOlive}{rgb}{0.1047,0.2412,0.0064}


\usepackage{fancyvrb}

\usepackage{mathptmx,helvet}
\usepackage[T1]{fontenc}
\usepackage{textcomp}


\usepackage[
            pdftex=true,
            bookmarks=true,        
            a4paper=true,
            pdftitle={Written with GAPDoc},
            pdfcreator={LaTeX with hyperref package / GAPDoc},
            colorlinks=true,
            backref=page,
            breaklinks=true,
            linkcolor=linkColor,
            citecolor=citeColor,
            filecolor=fileColor,
            urlcolor=urlColor,
            pdfpagemode={UseNone}, 
           ]{hyperref}

\newcommand{\maintitlesize}{\fontsize{50}{55}\selectfont}

% write page numbers to a .pnr log file for online help
\newwrite\pagenrlog
\immediate\openout\pagenrlog =\jobname.pnr
\immediate\write\pagenrlog{PAGENRS := [}
\newcommand{\logpage}[1]{\protect\write\pagenrlog{#1, \thepage,}}
%% were never documented, give conflicts with some additional packages

\newcommand{\GAP}{\textsf{GAP}}

%% nicer description environments, allows long labels
\usepackage{enumitem}
\setdescription{style=nextline}

%% depth of toc
\setcounter{tocdepth}{1}





%% command for ColorPrompt style examples
\newcommand{\gapprompt}[1]{\color{promptColor}{\bfseries #1}}
\newcommand{\gapbrkprompt}[1]{\color{brkpromptColor}{\bfseries #1}}
\newcommand{\gapinput}[1]{\color{gapinputColor}{#1}}


\begin{document}

\logpage{[ 0, 0, 0 ]}
\begin{titlepage}
\mbox{}\vfill

\begin{center}{\maintitlesize \textbf{ IdPGroup \\
\mbox{}}}\\
\vfill

\hypersetup{pdftitle={ IdPGroup }}
\markright{\scriptsize \mbox{}\hfill  IdPGroup  \hfill\mbox{}}
{\Huge \textbf{ Identification functions for different p\texttt{\symbol{45}}groups \\
\mbox{}}}\\
\vfill

{\Huge  1.0 \mbox{}}\\[1cm]
{ 11 August 2026 \mbox{}}\\[1cm]
\mbox{}\\[2cm]
{\Large \textbf{\strut  Henrik Schanze     \strut\mbox{}}}\\
\hypersetup{pdfauthor={ Henrik Schanze    }}
\end{center}\vfill

\mbox{}\\
{\mbox{}\\
\small \noindent \textbf{ Henrik Schanze    }  Email: \href{mailto://henrik.schanze@tu-braunschweig.de} {\texttt{henrik.schanze@tu\texttt{\symbol{45}}braunschweig.de}}\\
  Homepage: \href{https://www.iaa.tu-bs.de/henschan} {\texttt{https://www.iaa.tu\texttt{\symbol{45}}bs.de/henschan}}\\
  Address: \begin{minipage}[t]{8cm}\noindent
 Universitaetsplatz 2\\
 \end{minipage}
}\\
\end{titlepage}

\newpage\setcounter{page}{2}
\newpage

\def\contentsname{Contents\logpage{[ 0, 0, 1 ]}}

\tableofcontents
\newpage

     
\chapter{\textcolor{Chapter }{Introduction}}\label{Chapter_Introduction}
\logpage{[ 1, 0, 0 ]}
\hyperdef{L}{X7DFB63A97E67C0A1}{}
{
  

 The purpose of the \textsf{IdPGroup} package is to implement identification functions for different $p$\texttt{\symbol{45}}groups and thereby extend the functionality of the \textsf{SmallGrps} library. Currently this includes an identification function for groups of
order 512. 

 Details of the implementation and functionality of the identification
functions are dicussed in section \ref{Section_Identification}. 

 Section \ref{Section_Load} introduces the definitions of (strong) siblings and includes functions to load
all such sets of groups of certain order. 

 Functions that were used to calculate these lists of groups and which can be
used to calculate them for different orders can be found in section \ref{Section_Calculate}. 

 Finally section \ref{Section_Build} describes all funcions that were used to build the search tree in the
descirbed identification routine and which can be used to build other such
trees. }

   
\chapter{\textcolor{Chapter }{Functionality}}\label{Chapter_Functionality}
\logpage{[ 2, 0, 0 ]}
\hyperdef{L}{X87F1120883F5B4D0}{}
{
  

 This chapter lists all functions contained in the \textsf{IdPGroup} package. 
\section{\textcolor{Chapter }{Identifying groups via a search tree}}\label{Section_Identification}
\logpage{[ 2, 1, 0 ]}
\hyperdef{L}{X82BA59FA865D19E5}{}
{
  

 This section briefly describes the functions and variables of this package
associated to the identification of an arbitrary group of order 512 via a
search tree. By 'identify' a group $G$ we mean that we want to find \emph{id} such that \texttt{SmallGroup(512, id)} is isomorphic to $G$. This is done by calculating invariants and comparing them with the ones of
the remaining candidates. This process is organized in the tree \texttt{ID{\textunderscore}GROUP{\textunderscore}512{\textunderscore}TREE} (\ref{IDuScoreGROUPuScore512uScoreTREE}). 

 The used tree is implemented as nested records with the following components: 
\begin{itemize}
\item  "fp": a list of integers, each being the encoded fingerprint of groups
resulting from the specific test performed at the current level of the
identification routine. 
\item  "next": a list comprised of the resprective subbranches corresponding to the
fingerprints in "fp". This is either a list of ids or a branch implemented as
a record itself. 
\end{itemize}
 Depending on the depth in the tree each record also has one of the following
two components: 
\begin{itemize}
\item  "test": a string that indicates which test was used via \texttt{TestOnGroup} (\ref{TestOnGroup}) to create the listed fingerprints. 
\item  "desc": a list of integers that encode a description of the performed
powermap\texttt{\symbol{45}}test 
\end{itemize}
 

 The first levels of \texttt{ID{\textunderscore}GROUP{\textunderscore}512{\textunderscore}TREE} (\ref{IDuScoreGROUPuScore512uScoreTREE}) try to distinguish groups by the following invariants: 
\begin{itemize}
\item  the rank of the group, 
\item  the collected orders of all elements of the group, 
\item  the numbers of conjugacy classes of each length and representative order (the
order of the elements in the class) of the group, 
\item  the isomorphism type of the parent of the group, i.e. the quotient of the
group and the smallest non\texttt{\symbol{45}}trivial element of the
lower\texttt{\symbol{45}}$p$\texttt{\symbol{45}}central series, 
\item  the abelian invarants, i.e. the isomorphism type of the commutator factor
group of the group $G$ and 
\item  the cycle structure of the permutations induced on all clusters of conjugacy
classes of same length and representative order by the power map $x\mapsto x^3$. 
\end{itemize}
 

 If necessary these invariants are followed by invaraints resulting from power
maps evaluated on partitions of $G$. The initial partition is made of the conjugacy classes of $G$ grouped such that conjugacy classes of the same length and same representative
order belong to the same partitioning subset. The collected result of a power
map applied to a subset forms a invaraint used to distinguish groups in these
levels of the search tree. The describing parameters (e.g. which subset, which
power map) is determined by the content of the component "desc". After the
evaluation, the partition is refined according to the respective result. 

 If the preceding invaraints do not lead to a single group, the following (more
time consuming) invaraints are calculated. 
\begin{itemize}
\item  The isomorphism types of the non\texttt{\symbol{45}}trivial factors formed by
the group and elements of the smallest non\texttt{\symbol{45}}trivial element
of the lower\texttt{\symbol{45}}p\texttt{\symbol{45}}central series, 
\item  the isomorphism types of the non\texttt{\symbol{45}}trivial factors formed by
the group and elements of the center and 
\item  the isomorphism types of all maximal subgroups. 
\end{itemize}
 

 Remaining are 139 pairs and 1 triplet of groups of unidentified groups. All
these groups are distinguished by a random isomorphism test the same way as
groups of order 256 are finally distinguished in the \textsf{SmallGrps} library. 

\subsection{\textcolor{Chapter }{InfoId512}}
\logpage{[ 2, 1, 1 ]}\nobreak
\hyperdef{L}{X7E229FC579EB31F7}{}
{\noindent\textcolor{FuncColor}{$\triangleright$\enspace\texttt{InfoId512\index{InfoId512@\texttt{InfoId512}}
\label{InfoId512}
}\hfill{\scriptsize (info class)}}\\


 Info class for the identification routine of groups of order 512. }

 

\subsection{\textcolor{Chapter }{IdGroup512}}
\logpage{[ 2, 1, 2 ]}\nobreak
\hyperdef{L}{X81731B9A798C5985}{}
{\noindent\textcolor{FuncColor}{$\triangleright$\enspace\texttt{IdGroup512({\mdseries\slshape G[, opt]})\index{IdGroup512@\texttt{IdGroup512}}
\label{IdGroup512}
}\hfill{\scriptsize (function)}}\\
\textbf{\indent Returns:\ }
The id of the group of order 512 that is isomorphic to \mbox{\texttt{\mdseries\slshape G}}. 



 Identifies the group \mbox{\texttt{\mdseries\slshape G}} of order 512. This function accepts a record \mbox{\texttt{\mdseries\slshape opt}} of options with the following content: 
\begin{itemize}
\item  \mbox{\texttt{\mdseries\slshape tree}}: A search tree that was built using the invariants resulting from $tests$. 
\item  \mbox{\texttt{\mdseries\slshape anupq}}: A boolean that determines if $IsPqIsomorphicPGroup$ from the \textsf{Anupq}\texttt{\symbol{45}}Package should be used to find the to \mbox{\texttt{\mdseries\slshape G}} isomorphic group among the remaining candidates. 
\item  \mbox{\texttt{\mdseries\slshape path}}: A boolean that determines if the path through the tree shall be recorded
while identification and returned as well. 
\end{itemize}
 The default values of \mbox{\texttt{\mdseries\slshape opt}} is \texttt{rec(anupq}:=\texttt{false}, \texttt{tree:=ID{\textunderscore}GROUP{\textunderscore}512{\textunderscore}TREE,
path:=} \texttt{false} \texttt{)}. 
\begin{Verbatim}[commandchars=!@|,fontsize=\small,frame=single,label=Example]
  !gapprompt@gap>| !gapinput@G := Image(IsomorphismPermGroup(SmallGroup(512, 1984)));;|
  !gapprompt@gap>| !gapinput@IdGroup512(G);|
  1984
  !gapprompt@gap>| !gapinput@IdGroup512(G, rec(path := true));|
  [ 1984, rec( fp := [ 2, 8217, 5321 ], pos := [ 2, 188, 1 ],
          tests := [ "Rank", "ElmOrds", [ 40511009 ] ] ) ]
  !gapprompt@gap>| !gapinput@H := Image(IsomorphismPermGroup(SmallGroup(512, 449)));;|
  !gapprompt@gap>| !gapinput@IdGroup512(H, rec(anupq := false));|
  449
  !gapprompt@gap>| !gapinput@IdGroup512(H, rec(anupq := true));|
  449
\end{Verbatim}
 }

 

\subsection{\textcolor{Chapter }{ID{\textunderscore}GROUP{\textunderscore}512{\textunderscore}TREE}}
\logpage{[ 2, 1, 3 ]}\nobreak
\hyperdef{L}{X7E6424F37C14208E}{}
{\noindent\textcolor{FuncColor}{$\triangleright$\enspace\texttt{ID{\textunderscore}GROUP{\textunderscore}512{\textunderscore}TREE\index{ID{\textunderscore}GROUP{\textunderscore}512{\textunderscore}TREE@\texttt{ID{\textunderscore}}\-\texttt{G}\-\texttt{R}\-\texttt{O}\-\texttt{U}\-\texttt{P{\textunderscore}512{\textunderscore}}\-\texttt{TREE}}
\label{IDuScoreGROUPuScore512uScoreTREE}
}\hfill{\scriptsize (global variable)}}\\


 The search tree of groups of order 512 that can determine the vast majority of
groups. At the beginning only a small part is loaded. Other branches are
loaded if necessary. }

 

\subsection{\textcolor{Chapter }{LoadTotalTree}}
\logpage{[ 2, 1, 4 ]}\nobreak
\hyperdef{L}{X7BA3357E86269C3B}{}
{\noindent\textcolor{FuncColor}{$\triangleright$\enspace\texttt{LoadTotalTree({\mdseries\slshape })\index{LoadTotalTree@\texttt{LoadTotalTree}}
\label{LoadTotalTree}
}\hfill{\scriptsize (function)}}\\


 Loads the entire search tree \texttt{ID{\textunderscore}GROUP{\textunderscore}512{\textunderscore}TREE} (\ref{IDuScoreGROUPuScore512uScoreTREE}). }

 }

 
\section{\textcolor{Chapter }{Loading (strong) siblings}}\label{Section_Load}
\logpage{[ 2, 2, 0 ]}
\hyperdef{L}{X7EDB790F85B7753E}{}
{
  

 This section describes the functions used to load precalculated lists of
(strong) siblings, i.e. groups that fulfil the following definitions: 
\begin{itemize}
\item  \emph{subgroup\texttt{\symbol{45}}equivalent groups} are groups with a bijection between the sets of proper subgroups s.t. preimage
and image are isomorphic, 
\item  \emph{factor\texttt{\symbol{45}}equivalent groups} are groups with a bijection between the sets of proper factors s.t. preimage
and image are isomorphic, 
\item  \emph{siblings} are subgroup\texttt{\symbol{45}} and factor\texttt{\symbol{45}}equivalent
groups, where the bijection of proper subgroups preservs conjugacy and maps
the Frattini subgroup, lower and upper central series onto the respective
counterparts, 
\item  \emph{strong siblings} are siblings that are also Brauer Pairs. 
\end{itemize}
 All (strong) siblings of orders dividing 512 together with some trivial cases
are implemented in this package and are callable by the following functions. 

\subsection{\textcolor{Chapter }{Siblings}}
\logpage{[ 2, 2, 1 ]}\nobreak
\hyperdef{L}{X878A3FFE837F9EF6}{}
{\noindent\textcolor{FuncColor}{$\triangleright$\enspace\texttt{Siblings({\mdseries\slshape order, justids[, indices]})\index{Siblings@\texttt{Siblings}}
\label{Siblings}
}\hfill{\scriptsize (function)}}\\
\textbf{\indent Returns:\ }
A list of lists of siblings of size \mbox{\texttt{\mdseries\slshape order}}. The parameter \mbox{\texttt{\mdseries\slshape justids}} controls if just the ids or the groups (default) shall be returned. 



 Loads and returns a list of siblings of size \mbox{\texttt{\mdseries\slshape order}} if available. Currently available sizes are: 
\begin{itemize}
\item  all orders of the form $p\,q$ for $p$ and $q$ primes, 
\item  all powers of 2 smaller than $2^9$ 
\item  all powers of 3 smaller than $3^5$ 
\item  all powers of a prime smaller than $p^4$ and 
\item  all orders smaller than 256 except 60, 120, 160, 168, 180, 189, 192 and 240. 
\end{itemize}
 If no list is available for the provided order \texttt{fail} is returned. If a list \mbox{\texttt{\mdseries\slshape indices}} is given, only the respective entries are returned. 
\begin{Verbatim}[commandchars=!@|,fontsize=\small,frame=single,label=Example]
  !gapprompt@gap>| !gapinput@Siblings(101);|
  [ ]
  !gapprompt@gap>| !gapinput@Siblings(128);|
  [ [ <pc group of size 128 with 7 generators>,
      <pc group of size 128 with 7 generators> ],
    [ <pc group of size 128 with 7 generators>,
      <pc group of size 128 with 7 generators> ],
    [ <pc group of size 128 with 7 generators>,
      <pc group of size 128 with 7 generators> ] ]
  !gapprompt@gap>| !gapinput@Siblings(128, true);|
  [ [ 1317, 1322 ], [ 1327, 1329 ], [ 1597, 1598 ] ]
  !gapprompt@gap>| !gapinput@Siblings(256, true, [1..3]);|
  [ [ 227, 228 ], [ 258, 260, 261 ], [ 845, 846 ] ]
\end{Verbatim}
 }

 

\subsection{\textcolor{Chapter }{StrongSiblings}}
\logpage{[ 2, 2, 2 ]}\nobreak
\hyperdef{L}{X843063807B06E65B}{}
{\noindent\textcolor{FuncColor}{$\triangleright$\enspace\texttt{StrongSiblings({\mdseries\slshape order, justids[, indices]})\index{StrongSiblings@\texttt{StrongSiblings}}
\label{StrongSiblings}
}\hfill{\scriptsize (function)}}\\
\textbf{\indent Returns:\ }
A list of lists of strong siblings of \mbox{\texttt{\mdseries\slshape order}}. 



 Loads and returns a list of strong siblings of \mbox{\texttt{\mdseries\slshape order}} if available. The parameter \mbox{\texttt{\mdseries\slshape justids}} controls if just the ids or the groups (default) shall be returned. Currently
available sizes are: 
\begin{itemize}
\item  all orders of the form $p\,q$ for $p$ and $q$ primes, 
\item  all powers of 2 smaller than $2^9$ 
\item  all powers of 3 smaller than $3^5$ 
\item  all powers of a prime smaller than $p^4$ and 
\item  all orders smaller than 256 except 60, 120, 160, 168, 180, 189, 192 and 240. 
\end{itemize}
 If no list is available for the provided order \texttt{fail} is returned. If a list \mbox{\texttt{\mdseries\slshape indices}} is given, only the respective entries are returned. 
\begin{Verbatim}[commandchars=!@|,fontsize=\small,frame=single,label=Example]
  !gapprompt@gap>| !gapinput@StrongSiblings(256, true);|
  [ [ 1734, 1735 ], [ 1736, 1737 ], [ 1739, 1740 ], [ 1741, 1742 ],
    [ 3678, 3679 ], [ 4154, 4157 ], [ 4155, 4158 ], [ 4156, 4159 ] ]
  !gapprompt@gap>| !gapinput@StrongSiblings(256, false, [1..3]);|
  [ [ <pc group of size 256 with 8 generators>,
      <pc group of size 256 with 8 generators> ],
    [ <pc group of size 256 with 8 generators>,
      <pc group of size 256 with 8 generators> ],
    [ <pc group of size 256 with 8 generators>,
      <pc group of size 256 with 8 generators> ] ]
  !gapprompt@gap>| !gapinput@StrongSiblings(256, true, [1..3]);|
  [ [ 1734, 1735 ], [ 1736, 1737 ], [ 1739, 1740 ] ]
\end{Verbatim}
 }

 }

 
\section{\textcolor{Chapter }{Functions for calculating (strong) siblings}}\label{Section_Calculate}
\logpage{[ 2, 3, 0 ]}
\hyperdef{L}{X7AA188507AFCEC21}{}
{
  

 This section describes all functions for calculating (strong) siblings of a
given order. 

\subsection{\textcolor{Chapter }{InfoSiblings}}
\logpage{[ 2, 3, 1 ]}\nobreak
\hyperdef{L}{X87BC15127AE5D98F}{}
{\noindent\textcolor{FuncColor}{$\triangleright$\enspace\texttt{InfoSiblings\index{InfoSiblings@\texttt{InfoSiblings}}
\label{InfoSiblings}
}\hfill{\scriptsize (info class)}}\\


 Info class for the calculation of (strong) siblings. }

 

\subsection{\textcolor{Chapter }{CalculateSiblings}}
\logpage{[ 2, 3, 2 ]}\nobreak
\hyperdef{L}{X7A81D24885DD46A1}{}
{\noindent\textcolor{FuncColor}{$\triangleright$\enspace\texttt{CalculateSiblings({\mdseries\slshape order})\index{CalculateSiblings@\texttt{CalculateSiblings}}
\label{CalculateSiblings}
}\hfill{\scriptsize (function)}}\\
\textbf{\indent Returns:\ }
A list of lists of ids of groups of size \mbox{\texttt{\mdseries\slshape order}} that are siblings. 



 Determines all siblings among the groups of size \mbox{\texttt{\mdseries\slshape order}} by applying the tests \texttt{"MaxSubId"}, \texttt{"SpSubId"}, \texttt{"SubId"}, \texttt{"MaxFactId"} and \texttt{"FactId"} iterativly. The tests \texttt{"MaxSubId"} and \texttt{"MaxFactId"} are not necessary, but are efficient and eliminate some sets before applying
the more costly tests. 
\begin{Verbatim}[commandchars=!@|,fontsize=\small,frame=single,label=Example]
  !gapprompt@gap>| !gapinput@CalculateSiblings(128);|
  [ [ 1317, 1322 ], [ 1327, 1329 ], [ 1597, 1598 ] ]
\end{Verbatim}
 }

 

\subsection{\textcolor{Chapter }{CalculateStrongSiblings}}
\logpage{[ 2, 3, 3 ]}\nobreak
\hyperdef{L}{X7E2FC2447FF1E8FE}{}
{\noindent\textcolor{FuncColor}{$\triangleright$\enspace\texttt{CalculateStrongSiblings({\mdseries\slshape order})\index{CalculateStrongSiblings@\texttt{CalculateStrongSiblings}}
\label{CalculateStrongSiblings}
}\hfill{\scriptsize (function)}}\\
\textbf{\indent Returns:\ }
A list of lists of ids of groups of size \mbox{\texttt{\mdseries\slshape order}} that are strong siblings. 



 First calculats all siblings of size \mbox{\texttt{\mdseries\slshape order}} via \texttt{CalculateSiblings} (\ref{CalculateSiblings}) and then determines all strong siblings among them. 
\begin{Verbatim}[commandchars=!@|,fontsize=\small,frame=single,label=Example]
  !gapprompt@gap>| !gapinput@CalculateStrongSiblings(256);|
  [ [ 1734, 1735 ], [ 1736, 1737 ], [ 1739, 1740 ], [ 1741, 1742 ], [ 3678, 3679 ],
    [ 4154, 4157 ], [ 4155, 4158 ], [ 4156, 4159 ] ]
\end{Verbatim}
 }

 }

 
\section{\textcolor{Chapter }{Functions for building a search tree}}\label{Section_Build}
\logpage{[ 2, 4, 0 ]}
\hyperdef{L}{X7E8276BF79FDD898}{}
{
  

 This section describes all functions and variables connected to building a
search tree that can be used with \texttt{IdGroup512} (\ref{IdGroup512}) to identify groups via implemented invariants. 

 In short this is done be initializing a tree via calling \texttt{TestOnClusterList} (\ref{TestOnClusterList}) with a selected test on a list of (lists of) ids of groups of given order and
then iteratnig the following three steps 
\begin{itemize}
\item  Collecting the clusters of the curent tree via \texttt{ClusterList} (\ref{ClusterList}), 
\item  Performing a choosen test on these clusters via \texttt{TestOnClusterList} (\ref{TestOnClusterList}), 
\item  Extending the tree with these test results via \texttt{AddClusterTestResults} (\ref{AddClusterTestResults}). 
\end{itemize}
 A very effective set of invaraints can be obtained by calling \texttt{CocTestOnClusterList} (\ref{CocTestOnClusterList}) in the second step. This should only be done when the groups of each provided
cluster yields the same value when the test \texttt{CocList} is applied. Furthermore these invaraints should be used for as few as possible
clusters since \texttt{CocTestOnClusterList} (\ref{CocTestOnClusterList}) includes a search for suitable parameters than can take quite some time. 

\subsection{\textcolor{Chapter }{ClusterList}}
\logpage{[ 2, 4, 1 ]}\nobreak
\hyperdef{L}{X7812E18678076D7B}{}
{\noindent\textcolor{FuncColor}{$\triangleright$\enspace\texttt{ClusterList({\mdseries\slshape tree})\index{ClusterList@\texttt{ClusterList}}
\label{ClusterList}
}\hfill{\scriptsize (function)}}\\
\textbf{\indent Returns:\ }
A list of all clusters of groups in \mbox{\texttt{\mdseries\slshape tree}} (i.e. the leafs of \mbox{\texttt{\mdseries\slshape tree}}). 



 Iterates through \mbox{\texttt{\mdseries\slshape tree}} and collects information about each leaf as a record. This includes the size
of the leaf, the path through \mbox{\texttt{\mdseries\slshape tree}} leading to the leaf, the tests performed including descriptions of performed
powermap\texttt{\symbol{45}}tests and the ids of the groups contained in the
leaf. 
\begin{Verbatim}[commandchars=!@|,fontsize=\small,frame=single,label=Example]
  !gapprompt@gap>| !gapinput@tree := rec( fp := [ 63049, 75721, 4310 ],  next := [ [ 3, 8 ], [ 4 ],|
  !gapprompt@gap>| !gapinput@           rec( fp := [ 11744, 86646 ], next := [ [ 5 ], [ 6, 7 ] ],|
  !gapprompt@gap>| !gapinput@               test := "MaxSubId" ) ], test := "ElmOrds" );;|
  !gapprompt@gap>| !gapinput@ClusterList(tree);|
  [ rec( descl := [  ], ids := [ 3, 8 ], path := [ 63049 ], size := 2,
         tests := [ "ElmOrds" ] ),
    rec( descl := [  ], ids := [ 4 ], path := [ 75721 ], size := 1,
         tests := [ "ElmOrds" ] ),
    rec( descl := [  ], ids := [ 5 ], path := [ 4310, 11744 ], size := 1,
         tests := [ "ElmOrds", "MaxSubId" ] ),
    rec( descl := [  ], ids := [ 6, 7 ], path := [ 4310, 86646 ], size := 2,
         tests := [ "ElmOrds", "MaxSubId" ] ) ]
\end{Verbatim}
 }

 

\subsection{\textcolor{Chapter }{TestOnGroup}}
\logpage{[ 2, 4, 2 ]}\nobreak
\hyperdef{L}{X7C14B8DF7A4C0115}{}
{\noindent\textcolor{FuncColor}{$\triangleright$\enspace\texttt{TestOnGroup({\mdseries\slshape G, test[, opt]})\index{TestOnGroup@\texttt{TestOnGroup}}
\label{TestOnGroup}
}\hfill{\scriptsize (function)}}\\
\textbf{\indent Returns:\ }
The result of \mbox{\texttt{\mdseries\slshape test}} applied to \mbox{\texttt{\mdseries\slshape G}}, if desried encoded as an integer. 



 Performs the selected \mbox{\texttt{\mdseries\slshape test}} on the group \mbox{\texttt{\mdseries\slshape G}}. The following choices for \mbox{\texttt{\mdseries\slshape test}} are implemented: 
\begin{itemize}
\item  \mbox{\texttt{\mdseries\slshape "ElmOrds"}}: Collects the orders of all elements of \mbox{\texttt{\mdseries\slshape G}}. 
\item  \mbox{\texttt{\mdseries\slshape "Center"}}: Calculates the abelian invariants of the center of \mbox{\texttt{\mdseries\slshape G}}. 
\item  \mbox{\texttt{\mdseries\slshape "Rank"}}: Calculates the rank of \mbox{\texttt{\mdseries\slshape G}}. 
\item  \mbox{\texttt{\mdseries\slshape "AbelInv"}}: Computes the order of each term of the derived series and the abelian
Invariants of the derived factors of \mbox{\texttt{\mdseries\slshape G}}. 
\item  \mbox{\texttt{\mdseries\slshape "UCSId"}}: Identifies the proper subgroups of the upper central series of \mbox{\texttt{\mdseries\slshape G}} via \texttt{IdGroup}. 
\item  \mbox{\texttt{\mdseries\slshape "LCSId"}}: Identifies the proper subgroups of the lower central series of \mbox{\texttt{\mdseries\slshape G}} via \texttt{IdGroup}. 
\item  \mbox{\texttt{\mdseries\slshape "PCSId"}}: Identifies the proper subgroups of the p\texttt{\symbol{45}}central series
of \mbox{\texttt{\mdseries\slshape G}} via \texttt{IdGroup}. 
\item  \mbox{\texttt{\mdseries\slshape "SpSubId"}}: Computes the Frattini subgroup, the upper and lower central series of \mbox{\texttt{\mdseries\slshape G}} and identifies all terms if possible via \texttt{IdGroup}. 
\item  \mbox{\texttt{\mdseries\slshape "MaxSubId"}}: Collects the sizes and isomorphism types of all maximal subgroups of \mbox{\texttt{\mdseries\slshape G}} via \texttt{IdGroup}. 
\item  \mbox{\texttt{\mdseries\slshape "SubId"}}: Collects the sizes and isormorphism types of conjugacy classes of all proper
subgoups via \texttt{IdGroup}. 
\item  \mbox{\texttt{\mdseries\slshape "MaxFactId"}}: Collects the sizes and isomorphism types of all maximal factors of \mbox{\texttt{\mdseries\slshape G}} via \texttt{IdGroup}. 
\item  \mbox{\texttt{\mdseries\slshape "CentQuot"}}: Collects the sizes and isomorphism types of all quotients of \mbox{\texttt{\mdseries\slshape G}} with non\texttt{\symbol{45}}trivial elements of the center of \mbox{\texttt{\mdseries\slshape G}} via \texttt{IdGroup}. 
\item  \mbox{\texttt{\mdseries\slshape "LCSFactId"}}: Collects the isomorphism types of all quotients of \mbox{\texttt{\mdseries\slshape G}} with with cyclic subgroups generated by by non\texttt{\symbol{45}}trivial
elements of the last non\texttt{\symbol{45}}trivial term of the
p\texttt{\symbol{45}}central series. 
\item  \mbox{\texttt{\mdseries\slshape "FactId"}}: Identifies all proper factors of \mbox{\texttt{\mdseries\slshape G}} via \texttt{IdGroup}. 
\item  \mbox{\texttt{\mdseries\slshape "CocList"}}: Collects the number of conjugacy classes of a \mbox{\texttt{\mdseries\slshape G}} of same length and same representative order (the order of a representative). 
\item  \mbox{\texttt{\mdseries\slshape "PMCycSt"}}: Determines the cycle\texttt{\symbol{45}}structure invariant for the power
map $x\mapsto x^3$ on the clusters of conjugacy classes via \texttt{PowerMapCocCycSt}. 
\item  \mbox{\texttt{\mdseries\slshape "OutId"}}: Tries to identify the groups of outer automorphisms of \mbox{\texttt{\mdseries\slshape G}}, either via \texttt{IdGroup} (if available) or via its size and \texttt{TestOnGroup} with \mbox{\texttt{\mdseries\slshape "ElmOrds"}} and \mbox{\texttt{\mdseries\slshape "CocList"}} (\mbox{\texttt{\mdseries\slshape "MaxSubId"}} and \mbox{\texttt{\mdseries\slshape "CentQuot"}} can also be used if feasible). 
\item  \mbox{\texttt{\mdseries\slshape "LattEqv"}}: Calculates all conjugacy classes of proper subgoups of \mbox{\texttt{\mdseries\slshape G}}, stores their size, identifies the the isomorphism types of each class and
determines the isomorphism types and conjugacy classes in \mbox{\texttt{\mdseries\slshape G}} of its maximal subgroups and minimal supergroups. 
\end{itemize}
 This function accepts a record \mbox{\texttt{\mdseries\slshape opt}} of options with the single content \mbox{\texttt{\mdseries\slshape opt.pack}}, a boolean that determines if the result of a test should be encoded as an
integer. The default value of \mbox{\texttt{\mdseries\slshape opt}} is $rec(pack:=$\texttt{true}$)$. 
\begin{Verbatim}[commandchars=!@|,fontsize=\small,frame=single,label=Example]
  !gapprompt@gap>| !gapinput@G := SmallGroup(512, 1405);;|
  !gapprompt@gap>| !gapinput@TestOnGroup(G, "ElmOrds");|
  85416
  !gapprompt@gap>| !gapinput@TestOnGroup(G, "ElmOrds", rec(pack := false));|
  [ [ 1, 1 ], [ 2, 7 ], [ 4, 24 ], [ 8, 288 ], [ 16, 64 ], [ 32, 128 ] ]
\end{Verbatim}
 }

 

\subsection{\textcolor{Chapter }{TestOnClusterList}}
\logpage{[ 2, 4, 3 ]}\nobreak
\hyperdef{L}{X7BC4B6397C17B535}{}
{\noindent\textcolor{FuncColor}{$\triangleright$\enspace\texttt{TestOnClusterList({\mdseries\slshape CL, test[, opt]})\index{TestOnClusterList@\texttt{TestOnClusterList}}
\label{TestOnClusterList}
}\hfill{\scriptsize (function)}}\\
\textbf{\indent Returns:\ }
Either a list of pairs each containg the path to the respective cluster (which
is a leaf of a search tree) and the subbranch resulting from the results of \mbox{\texttt{\mdseries\slshape test}}, or just the list of subbranches if \mbox{\texttt{\mdseries\slshape CL}} contains only lists of integers. If \mbox{\texttt{\mdseries\slshape CL}} consits of a single list of integers the single resulting subbranch is
returned on its own. 



 Performs the choosen test on the groups of each record in the \mbox{\texttt{\mdseries\slshape CL}}. The tests are performed by \texttt{TestOnGroup} (\ref{TestOnGroup}). \mbox{\texttt{\mdseries\slshape CL}} could either be a list of records produced by \texttt{ClusterList} (\ref{ClusterList}) or be a list of lists containing ids. In the former case the output is a list
of pairs containg the path of a cluster and the subbranch resulting from the
test. In the latter case only a list of subbranches is returned. If in the
latter case the returned list only one contains one subbranch, the subbranch
itself is returned. This is useful when building the first level of a search
tree. This function accepts a record \mbox{\texttt{\mdseries\slshape opt}} of options with the following content: 
\begin{itemize}
\item  \mbox{\texttt{\mdseries\slshape pack}}: A boolean that is handed to \texttt{TestOnGroup} (\ref{TestOnGroup}). 
\item  \mbox{\texttt{\mdseries\slshape ign1cl}}: A boolean that determines if clusters of size 1 should be ignored. 
\item  \mbox{\texttt{\mdseries\slshape order}}: An integer giving the order of the groups with ids in \mbox{\texttt{\mdseries\slshape CL}}. 
\item  \mbox{\texttt{\mdseries\slshape verbose}}: A boolean that determines if progress should be printed. 
\item  \mbox{\texttt{\mdseries\slshape notetest}}: A boolean that decides if the performed test should be noted in the returned
subbranches. If so subbranches, that do not split a cluster are omited. 
\end{itemize}
 The default value of \mbox{\texttt{\mdseries\slshape opt}} is $rec(pack$:=\texttt{true}, $ign1cl$:=\texttt{true}, $order$:=512, $verbose$:=\texttt{true}, $notetest$:=\texttt{true}$)$. 
\begin{Verbatim}[commandchars=!@|,fontsize=\small,frame=single,label=Example]
  !gapprompt@gap>| !gapinput@TestOnClusterList([[1..10]],"ElmOrds", rec(verbose:=false));|
  rec( fp := [ 88570, 4310, 63049, 75721 ],
       next := [ [ 1 ], [ 2, 5, 6, 7, 9, 10 ], [ 3, 8 ], [ 4 ] ],
       test := "ElmOrds")
  !gapprompt@gap>| !gapinput@TestOnClusterList([[1,2,5],[7,8,9]],"ElmOrds", rec(verbose:=false));|
  [ rec( fp := [ 88570, 4310 ], next := [ [ 1 ], [ 2, 5 ] ],
         test := "ElmOrds" ),
    rec( fp := [ 4310, 63049 ], next := [ [ 7, 9 ], [ 8 ] ],
         test := "ElmOrds") ]
  !gapprompt@gap>| !gapinput@tree := TestOnClusterList([[3..8]],"ElmOrds", rec(verbose:=false));;|
  !gapprompt@gap>| !gapinput@CL := ClusterList(tree);|
  [ rec( descl := [  ], ids := [ 3, 8 ], path := [ 63049 ], size := 2,
         tests := [ "ElmOrds" ]),
    rec( descl := [  ], ids := [ 4 ], path := [ 75721 ], size := 1,
         tests := [ "ElmOrds" ]),
    rec( descl := [  ], ids := [ 5, 6, 7 ], path := [ 4310 ], size := 3,
         tests := [ "ElmOrds" ]) ]
  !gapprompt@gap>| !gapinput@TestOnClusterList(CL, "MaxSubId", rec(verbose:=false, notetest := false));|
  [ [ [ 63049 ], rec( fp := [ 62603 ], next := [ [ 3, 8 ] ] ) ],
    [ [ 4310 ], rec( fp := [ 11744, 86646 ], next := [ [ 5 ], [ 6, 7 ] ] ) ] ]
\end{Verbatim}
 }

 

\subsection{\textcolor{Chapter }{AddClusterTestResults}}
\logpage{[ 2, 4, 4 ]}\nobreak
\hyperdef{L}{X85F0BFE88445D1B7}{}
{\noindent\textcolor{FuncColor}{$\triangleright$\enspace\texttt{AddClusterTestResults({\mdseries\slshape tree, TR})\index{AddClusterTestResults@\texttt{AddClusterTestResults}}
\label{AddClusterTestResults}
}\hfill{\scriptsize (function)}}\\
\textbf{\indent Returns:\ }
The extended \mbox{\texttt{\mdseries\slshape tree}} with the subbranches in \mbox{\texttt{\mdseries\slshape TR}}. 



 Extends \mbox{\texttt{\mdseries\slshape tree}} by the test results stored in \mbox{\texttt{\mdseries\slshape TR}}, a list of pairs each containing the path in \mbox{\texttt{\mdseries\slshape tree}} and a branch (a tree of depth one), extending the corresponding cluster of the
inital \mbox{\texttt{\mdseries\slshape tree}}. 
\begin{Verbatim}[commandchars=!@|,fontsize=\small,frame=single,label=Example]
  !gapprompt@gap>| !gapinput@tree := TestOnClusterList([[3..8]],"ElmOrds", rec(verbose:=false));|
  rec( fp := [ 63049, 75721, 4310 ], next := [ [ 3, 8 ], [ 4 ], [ 5, 6, 7 ] ],
       test := "ElmOrds")
  !gapprompt@gap>| !gapinput@TR:=TestOnClusterList(ClusterList(tree),"MaxSubId",rec(verbose:=false));|
  [ [ [ 4310 ], rec( fp := [ 11744, 86646 ], next := [ [ 5 ], [ 6, 7 ] ],
                     test := "MaxSubId" ) ] ]
  !gapprompt@gap>| !gapinput@AddClusterTestResults(tree,TR);|
  rec( fp := [ 63049, 75721, 4310 ],
       next := [ rec( fp := [ 62603 ], next := [ [ 3, 8 ] ] ), [ 4 ],
                 rec( fp := [ 11744, 86646 ], next := [ [ 5 ], [ 6, 7 ] ],
       test := "MaxSubId") ], test := "ElmOrds" )
\end{Verbatim}
 }

 

\subsection{\textcolor{Chapter }{BuildSearchTree}}
\logpage{[ 2, 4, 5 ]}\nobreak
\hyperdef{L}{X82CDC04E8588770B}{}
{\noindent\textcolor{FuncColor}{$\triangleright$\enspace\texttt{BuildSearchTree({\mdseries\slshape N, order, tests})\index{BuildSearchTree@\texttt{BuildSearchTree}}
\label{BuildSearchTree}
}\hfill{\scriptsize (function)}}\\
\textbf{\indent Returns:\ }
A search tree (implemented as a record). 



 Builds a search tree for the groups with ids in \mbox{\texttt{\mdseries\slshape N}} or, if \mbox{\texttt{\mdseries\slshape N}} is an integer, the first \mbox{\texttt{\mdseries\slshape N}} groups of size \mbox{\texttt{\mdseries\slshape order}}, using \texttt{TestOnClusterList} (\ref{TestOnClusterList}) with the tests from the list \mbox{\texttt{\mdseries\slshape tests}}. 
\begin{Verbatim}[commandchars=!@|,fontsize=\small,frame=single,label=Example]
  !gapprompt@gap>| !gapinput@BuildSearchTree(5,512,["ElmOrds", "MaxSubId"]);|
  rec( fp := [ 88570, 4310, 63049, 75721 ], next := [ [ 1 ],
       rec( fp := [ 76646, 11744 ], next := [ [ 2 ], [ 5 ] ],
       test := "MaxSubId" ), [ 3 ], [ 4 ] ], test := "ElmOrds" )
  !gapprompt@gap>| !gapinput@BuildSearchTree([1,2,3,5,7], 512, ["ElmOrds", "CocList"]);|
  rec( fp := [ 88570, 4310, 63049 ], next := [ [ 1 ],
       rec( fp := [ 96561, 32173 ], next := [ [ 2 ], [ 5, 7 ] ],
       test := "CocList" ), [ 3 ] ], test := "ElmOrds" )
\end{Verbatim}
 }

 

\subsection{\textcolor{Chapter }{CocTestOnClusterList}}
\logpage{[ 2, 4, 6 ]}\nobreak
\hyperdef{L}{X7EB7CB0B84963812}{}
{\noindent\textcolor{FuncColor}{$\triangleright$\enspace\texttt{CocTestOnClusterList({\mdseries\slshape CL[, opt]})\index{CocTestOnClusterList@\texttt{CocTestOnClusterList}}
\label{CocTestOnClusterList}
}\hfill{\scriptsize (function)}}\\
\textbf{\indent Returns:\ }
A list of pairs each containing the path of the respective cluster and the
subbranch resulting from \texttt{CocSplitRoutine}. 



 Tries to split each cluster in \mbox{\texttt{\mdseries\slshape CL}} via tests on the lists of conjugacy classes. \mbox{\texttt{\mdseries\slshape CL}} must be a cluster list produced by \texttt{ClusterList} (\ref{ClusterList}). This function accepts a record \mbox{\texttt{\mdseries\slshape opt}} of options with the following content: 
\begin{itemize}
\item  \mbox{\texttt{\mdseries\slshape pack}}: A boolean that determines if the result should be encoded as an intger. 
\item  \mbox{\texttt{\mdseries\slshape order}}: An integer giving the order of the groups with ids in \mbox{\texttt{\mdseries\slshape CL}}. 
\item  \mbox{\texttt{\mdseries\slshape iterate}}: An integer that determines the maximal number of iterations done while
searching parameters. 
\item  \mbox{\texttt{\mdseries\slshape verbose}}: An boolean that determines if progress should be printed. 
\end{itemize}
 The default value of \mbox{\texttt{\mdseries\slshape opt}} is \mbox{\texttt{\mdseries\slshape rec(pack}}:=\texttt{true}, \mbox{\texttt{\mdseries\slshape order}}:=512, \mbox{\texttt{\mdseries\slshape iterate}}:=1,\mbox{\texttt{\mdseries\slshape verbose}}:=\texttt{true}\mbox{\texttt{\mdseries\slshape )}}. 
\begin{Verbatim}[commandchars=!@|,fontsize=\small,frame=single,label=Example]
  !gapprompt@gap>| !gapinput@tree := TestOnClusterList([[3,5,6,7,8]],"CocList",rec(verbose:=false));;|
  !gapprompt@gap>| !gapinput@CL := ClusterList(tree);|
  [ rec( descl := [  ], ids := [ 3, 8 ], path := [ 337 ], size := 2,
         tests := [ "CocList" ] ),
    rec( descl := [  ], ids := [ 5, 6, 7 ], path := [ 32173 ], size := 3,
         tests := [ "CocList" ] ) ]
  !gapprompt@gap>| !gapinput@TR := CocTestOnClusterList(CL,rec(verbose:=false));;|
  !gapprompt@gap>| !gapinput@AddClusterTestResults(tree, TR);;|
  !gapprompt@gap>| !gapinput@ClusterList(tree);|
  [ rec( descl := [ [ 302004 ] ], ids := [ 3 ], path := [ 337, 1363 ], size := 1,
         tests := [ "CocList", "" ]),
    rec( descl := [ [ 302004 ] ], ids := [ 8 ], path := [ 337, 46760 ], size := 1,
         tests := [ "CocList", ""  ]),
    rec( descl := [ [ 104003 ] ], ids := [ 5 ], path := [ 32173, 1214 ],
         size := 1, tests := [ "CocList", ""  ] ),
    rec( descl := [ [ 104003 ], [ 40204005 ] ], ids := [ 6 ],
         path := [ 32173, 16, 20792 ], size := 1, tests := [ "CocList", ""  ] ),
    rec( descl := [ [ 104003 ], [ 40204005 ] ], ids := [ 7 ],
         path := [ 32173, 16, 46881 ], size := 1, tests := [ "CocList", ""  ] ) ]
\end{Verbatim}
 }

 

\subsection{\textcolor{Chapter }{ParametersRandPresentTest}}
\logpage{[ 2, 4, 7 ]}\nobreak
\hyperdef{L}{X7FBD2B9C83B1496A}{}
{\noindent\textcolor{FuncColor}{$\triangleright$\enspace\texttt{ParametersRandPresentTest({\mdseries\slshape CL, order, maxl})\index{ParametersRandPresentTest@\texttt{ParametersRandPresentTest}}
\label{ParametersRandPresentTest}
}\hfill{\scriptsize (function)}}\\
\textbf{\indent Returns:\ }
A list consisting of pairs of a list of ids and a record containing the
parameters for a random presentation test via \texttt{ID{\textunderscore}GROUP{\textunderscore}FUNCS[18]}. 



 Determines the needed parameters for the groups of size \mbox{\texttt{\mdseries\slshape order}} in the clusters \mbox{\texttt{\mdseries\slshape cl}} of \mbox{\texttt{\mdseries\slshape CL}}. Creates the clusters of conjugcy classes and alters them accordingly to the \mbox{\texttt{\mdseries\slshape cl.descl}}. Afterwards minimal generators for all groups and their position among the
clusters of conjugacy classes. If the positions do not coincide, the positions
of the generators of the first group are selected and the generators of the
other groups are selected from the same clusters. Otherwise the positions and
generators are sorted by the positions. Sequently the function tries to refine
the corresponding clusters to limit the choices to choose from and repeat the
process until either the number of refinments is equal or larger than \mbox{\texttt{\mdseries\slshape maxl}} or no more refinments could be found. 
\begin{Verbatim}[commandchars=!@|,fontsize=\small,frame=single,label=Example]
  !gapprompt@gap>| !gapinput@ParametersRandPresentTest([rec(descl:=[], ids:=[1367, 1371])], 512, 0);|
  [ [ [ 1367, 1371 ], rec( descl := [  ],
       fp := [ 39080, 80814 ], func := 18, pos := [ 8, 10 ] ) ] ]
  !gapprompt@gap>| !gapinput@ParametersRandPresentTest([rec(descl:=[], ids:=[1367, 1371])], 512, 1);|
  [ [ [ 1367, 1371 ], rec( descl := [ 106003 ],
       fp := [ 39080, 80814 ], func := 18, pos := [ 9, 11 ] ) ] ]
\end{Verbatim}
 }

 }

 }

 \def\indexname{Index\logpage{[ "Ind", 0, 0 ]}
\hyperdef{L}{X83A0356F839C696F}{}
}

\cleardoublepage
\phantomsection
\addcontentsline{toc}{chapter}{Index}


\printindex

\immediate\write\pagenrlog{["Ind", 0, 0], \arabic{page},}
\immediate\write\pagenrlog{["Ind", 0, 0], \arabic{page},}
\immediate\write\pagenrlog{["Ind", 0, 0], \arabic{page},}
\immediate\write\pagenrlog{["Ind", 0, 0], \arabic{page},}
\immediate\write\pagenrlog{["Ind", 0, 0], \arabic{page},}
\immediate\write\pagenrlog{["Ind", 0, 0], \arabic{page},}
\immediate\write\pagenrlog{["Ind", 0, 0], \arabic{page},}
\immediate\write\pagenrlog{["Ind", 0, 0], \arabic{page},}
\newpage
\immediate\write\pagenrlog{["End"], \arabic{page}];}
\immediate\closeout\pagenrlog
\end{document}
